\section{PRIMARY KEY dan UNIQUE}

\textbf{PRIMARY KEY} dan \textbf{UNIQUE} adalah dua jenis constraint yang menjamin keunikan data, namun memiliki perbedaan penting. Primary key mengidentifikasi setiap baris secara unik dan \textbf{tidak boleh NULL}; setiap tabel hanya boleh memiliki satu primary key (bisa terdiri dari beberapa kolom --- \textit{composite key}). UNIQUE juga menjamin keunikan nilai, namun \textbf{boleh berisi NULL} (satu nilai NULL per kolom) dan satu tabel boleh memiliki beberapa constraint UNIQUE \cite{mysqlDocs}.

\begin{table}[ht]
\centering
\begin{tabular}{|P{3.5cm}|P{5cm}|P{5cm}|}
\hline
\textbf{Aspek} & \textbf{PRIMARY KEY} & \textbf{UNIQUE} \\
\hline
Jumlah per tabel & Hanya 1 & Boleh banyak \\
\hline
Boleh NULL? & Tidak & Ya (1 NULL) \\
\hline
Index otomatis & Clustered index & Non-clustered index \\
\hline
Fungsi utama & Identitas baris & Mencegah duplikasi \\
\hline
Contoh & id, nim & email, username \\
\hline
\end{tabular}
\caption{Perbedaan PRIMARY KEY dan UNIQUE}
\end{table}

\begin{webcode}[language=SQL, caption={Penggunaan PRIMARY KEY dan UNIQUE}]
CREATE TABLE pengguna (
  id       INT AUTO_INCREMENT PRIMARY KEY,
  username VARCHAR(50) NOT NULL UNIQUE,
  email    VARCHAR(100) NOT NULL UNIQUE,
  nama     VARCHAR(100) NOT NULL
);

-- Composite primary key
CREATE TABLE pesanan_detail (
  pesanan_id INT,
  produk_id  INT,
  jumlah     INT NOT NULL DEFAULT 1,
  PRIMARY KEY (pesanan_id, produk_id)
);
\end{webcode}

